%%%% CAPÍTULO 3 - MATERIAL E MÉTODOS

\chapter{Materiais e Métodos}\label{cap:materialemetodos}

Este capítulo descreve os materiais e os métodos utilizados para construir e validar a plataforma proposta. A ênfase metodológica está no uso de MATLAB HDL Coder para gerar a lógica de estratégia, enquanto a infraestrutura de comunicação, \gls{orderbook} e testes permanece como base de execução.

\section{Materiais}\label{sec:materiais}

A plataforma alvo é a DE10-Nano, baseada em Intel Cyclone V SoC. A escolha dessa placa é adequada ao trabalho porque combina processador ARM, FPGA e interface HPS--FPGA em um ambiente de baixo custo. O ARM executa Linux embarcado, sockets e decodificação de mensagens; a FPGA executa os blocos que exigem paralelismo e comportamento temporal mais previsível.

O software de apoio foi desenvolvido em C++ e utiliza a biblioteca mFAST para codificação e decodificação de mensagens FAST. Essa camada é responsável por gerar dados sintéticos, receber mensagens, normalizar símbolos, converter preços para ponto fixo e interagir com a região de memória mapeada da FPGA. Essa escolha mantém no software as partes mais variáveis do protocolo de mercado e envia ao hardware apenas dados numéricos.

O hardware foi descrito em VHDL, com simulações também em VHDL. A infraestrutura manual inclui a ponte ARM--FPGA, o \gls{orderbook}, o empacotamento de respostas e a adaptação ao barramento Avalon-MM. A estratégia de trading, por sua vez, é escrita em MATLAB e convertida para VHDL por meio do HDL Coder. Essa divisão foi adotada para preservar controle explícito das interfaces críticas e, ao mesmo tempo, facilitar a criação de novas regras de decisão.

O fluxo de construção e validação utiliza Make, Docker, CMake, GHDL, MATLAB HDL Coder e Quartus Prime Lite com Platform Designer. Essas ferramentas são empregadas para reduzir diferenças de ambiente, executar testes C++ e VHDL, gerar o bloco HDL da estratégia, compilar o projeto FPGA e preparar binários para a DE10-Nano. Como o trabalho depende de ferramentas acadêmicas e dados sintéticos, não há necessidade de acesso a uma bolsa real, contrato de dados proprietário ou infraestrutura profissional de negociação.

\section{Métodos}\label{sec:metodo}

\subsection{Organização da plataforma}\label{subsec:definicaoPlataforma}

O método de desenvolvimento foi orientado pela separação entre infraestrutura e estratégia. A infraestrutura compreende recepção de dados, normalização, transporte por MMIO, \gls{orderbook}, integração Avalon-MM e testes. A estratégia é mantida em uma função MATLAB sintetizável, com entradas numéricas já derivadas do \gls{orderbook}. Assim, a troca de estratégia deve exigir alteração prioritariamente no código MATLAB, não na ponte de comunicação nem no \gls{orderbook}.

Essa organização também define a responsabilidade de cada domínio. O software lida com mensagens textuais ou binárias de formato flexível. O hardware lida com palavras de largura fixa, estado do livro e decisão. O MATLAB HDL Coder atua no ponto de maior interesse do usuário: a regra de trading. A consequência é um fluxo de trabalho mais simples para experimentar variações de estratégia, desde que sejam respeitadas as restrições de código sintetizável.

\subsection{Integração no Quartus e Platform Designer}\label{subsec:metodoQuartus}

Além da descrição RTL e da geração automática da estratégia, o trabalho exige uma etapa explícita de integração no Quartus Prime Lite, em especial no Platform Designer da Intel \cite{IntelPlatformDesigner}. É nessa ferramenta que o \gls{hps}, a ponte HPS--FPGA e o periférico customizado \codigo{hft\_trade\_engine\_avalon\_mm} são ligados em um sistema SoC coerente para a DE10-Nano.

A \autoref{fig:platformDesigner} mostra a integração utilizada no projeto. Nela, o HPS atua como bloco iniciador da comunicação e o motor HFT aparece como periférico Avalon-MM do tipo \textit{slave}. Essa ligação é importante porque materializa, em nível de ferramenta, o contrato discutido ao longo do texto: o software ARM publica eventos pela ponte, o periférico no FPGA consome esses eventos, atualiza o \gls{orderbook} e devolve respostas pela fila RX. Em outras palavras, a figura serve como evidência visual de como a arquitetura proposta foi efetivamente montada no ambiente Quartus.

\begin{figure}[htb]
\caption{Ligação entre HPS e o bloco \textit{HFT Trade Engine Avalon-MM} no Platform Designer}
\label{fig:platformDesigner}
\centering
\includegraphics[width=0.95\textwidth]{platform-designer}
\fonte{Captura de tela elaborada pelos autores no Quartus Prime Lite Platform Designer.}
\end{figure}

Essa etapa de integração complementa os artefatos em VHDL e os scripts de build. Embora a lógica de decisão, a ponte ARM--FPGA e o \gls{orderbook} estejam descritos nos arquivos-fonte, é o Platform Designer que define a conexão final com o HPS, o mapeamento Avalon-MM e a forma como o sistema sintetizado passa a existir como periférico acessível pelo software embarcado.

\subsection{Contrato de comunicação}\label{subsec:metodoContrato}

O contrato ARM--FPGA foi definido como uma interface de produtor e consumidor baseada em duas filas circulares: uma para eventos enviados pelo ARM e outra para respostas produzidas pela FPGA. Cada evento transporta sequência, símbolo normalizado, lado, preço em ponto fixo e quantidade. Cada resposta devolve sequência, ação, melhor compra, melhor venda, \textit{spread} e desbalanceamento.

O uso de filas evita que cada evento dependa de uma chamada síncrona por registradores individuais. Com isso, o ARM pode publicar eventos enquanto houver espaço e a FPGA pode consumi-los conforme sua disponibilidade. Para manter corretude, o produtor escreve o conteúdo completo antes de publicar o ponteiro, e o consumidor lê o conteúdo completo antes de liberar o espaço. Os detalhes de offsets e palavras individuais são apresentados nas tabelas e apêndices, pois pertencem mais ao contrato técnico do que à discussão metodológica principal.

\subsection{Fluxo de dados de mercado}\label{subsec:metodoFeed}

O fluxo de dados usa mensagens sintéticas codificadas em FAST e transmitidas por TCP. A escolha por um feed sintético permite repetir testes sem depender de dados reais ou infraestrutura de mercado. O receptor decodifica as mensagens, converte campos para representação numérica e, quando a FPGA está habilitada, envia o evento normalizado ao hardware.

Essa abordagem não pretende medir a latência de uma rede de bolsa real. O objetivo é exercitar o caminho de dados até a estratégia: mensagem de mercado, normalização, transporte ARM--FPGA, atualização do livro e retorno da decisão. Assim, a metodologia separa validação funcional do pipeline de qualquer afirmação quantitativa sobre redes financeiras reais.

\subsection{Order Book e estratégia}\label{subsec:metodoLivroEstrategia}

O \gls{orderbook} foi implementado em hardware com profundidade limitada. Ele mantém níveis de compra e venda por símbolo, atualiza preços e quantidades, rejeita estados cruzados e calcula os sinais necessários à estratégia. A profundidade limitada é coerente com o objetivo do protótipo: validar a integração e fornecer sinais de topo de livro, não reproduzir integralmente a complexidade de um livro de bolsa.

A estratégia foi escrita em MATLAB com operações compatíveis com HDL Coder. Ela recebe os sinais do \gls{orderbook} e retorna uma ação discreta: NOOP, BUY ou SELL. O papel do HDL Coder é transformar essa regra em lógica sintetizável, permitindo que o usuário trabalhe em um nível mais próximo da formulação algorítmica. Estratégias futuras podem alterar limiares ou combinar os sinais de outra forma, desde que mantenham tipos e operações adequados à síntese.

\subsection{Validação funcional}\label{subsec:metodoValidacao}

A validação foi dividida por camadas. O software é validado por testes de comunicação e empacotamento. A ponte ARM--FPGA é validada por simulações de fila, rajadas e contrapressão. O \gls{orderbook} é validado por casos de inserção, remoção, ordenação e rejeição de estados inválidos. A estratégia é validada por cenários que devem produzir NOOP, BUY e SELL. Por fim, o motor completo é validado pela passagem de eventos por todas as etapas.

O \autoref{quad:matrizValidacao} resume os principais riscos cobertos por essa validação. A matriz evita que o trabalho dependa apenas de uma demonstração manual em placa, pois cada parte essencial da arquitetura possui uma evidência funcional própria.

\begin{tabframed}[htb]
\caption{Matriz de validação funcional do protótipo}
\label{quad:matrizValidacao}
\begin{tabularx}{\textwidth}{lXX}
\toprule
\textbf{Artefato} & \textbf{Risco técnico mitigado} & \textbf{Evidência de validação} \\ \midrule
Interface ARM--FPGA & Divergência de contrato, fila cheia mal detectada ou avanço de ponteiro fora de ordem. & Teste C++ e testbenches VHDL da ponte. \\
Order Book & Livro fora de ordem, remoção incorreta ou aceitação de livro cruzado. & Testbench específico com casos de inserção, atualização e remoção. \\
Estratégia gerada & Ação incompatível com os sinais de mercado ou resposta perdida. & Testbench da estratégia e verificação de NOOP em estados inválidos. \\
Motor integrado & Falha de integração entre comunicação, livro e decisão. & Testbench de motor completo com escrita de eventos e leitura de respostas. \\
Integração Avalon-MM & Barramento travado ou transação sem retorno durante acesso do HPS. & Testbench do wrapper Avalon-MM. \\
Feed e receptor & Falha de codificação, decodificação ou normalização dos eventos. & Teste de fumaça software-only e execução funcional em placa. \\ \bottomrule
\end{tabularx}
\fonte{}
\end{tabframed}

\subsection{Método proposto para métricas experimentais}\label{subsec:metodoMedicoes}

A implementação atual fornece principalmente evidências funcionais. Para completar a avaliação quantitativa, o método prevê instrumentar o software e a FPGA. No software, devem ser registrados os instantes de recebimento do evento, publicação do quadro e leitura da resposta. No hardware, recomenda-se usar um contador de ciclos para medir o intervalo entre consumo do evento e publicação da resposta.

A vazão deve ser medida aumentando progressivamente a taxa do feed até ocorrer saturação, fila cheia ou perda de resposta.

\begin{tabframed}[htb]
\caption{Métricas experimentais planejadas e resultados obtidos}
\label{quad:metricasPlanejadas}
\begin{tabularx}{\textwidth}{lXXX}
\toprule
\textbf{Métrica} & \textbf{Como medir} & \textbf{Meta de referência} & \textbf{Resultado} \\ \midrule
Latência fim a fim (RTT) & Modo \codigo{fpga-sync}: mede RTT ARM$\to$FPGA$\to$ARM por mensagem com percentis. & Comparar com versão software-only. & \textbf{10.228~ns} médio; \textbf{10.310~ns} p99 (\autoref{tab:latenciaComparacao}) \\
Latência interna FPGA & Contador de ciclos do núcleo lógico. & Jitter interno inferior a $1~\mu$s. & \textbf{60~ns, jitter 0~ns} (\autoref{sec:benchmarkDesempenho}) \\
Vazão & Eventos por segundo até ocorrência de saturação. & Pelo menos 100.000 mensagens/s. & \textbf{102.381~msg/s} \\
Speedup (núcleo lógico) & Comparação do núcleo de decisão FPGA vs.\ C++ equivalente (sem comun.). & Aceleração mínima de 5 vezes. & \textbf{5,253\texttimes} (lógica isolada) \\ \bottomrule
\end{tabularx}
\fonte{}
\end{tabframed}

As metas apresentadas funcionaram como critérios de avaliação para a campanha experimental e, na versão atual do trabalho, já podem ser confrontadas com resultados efetivamente obtidos. Essa distinção continua importante para não confundir integração funcional com desempenho: o benchmark já comprova aceleração do núcleo lógico e vazão da ponte MMIO.
